\section{Results}

\subsection{Which encoders expose compatibility to a linear probe?}
The leader depends on the corpus (Fig.~\ref{fig:ranking}). SigLIP2-So400m reaches a native-probe CP-AUC of \resPolyvoreVTwoCleanSiglipTwoSoFourZeroZeromThreeEightFourNativeLogisticTestCpAuc{} on Polyvore. Qwen3-VL-Embedding-8B reaches \resIqonThreeZeroZeroZeroVTwoCleanQwenThreeVlEmbEightbNativeLogisticTestCpAuc{} on IQON, while PE-Core reaches \resPogVTwoCleanPeCoreLOneFourThreeThreeSixNativeLogisticTestCpAuc{} on POG. Qwen leads whole-outfit discrimination on IQON, whereas SigLIP2-So400m recovers the missing item more often there. ResNet-50 ranks last on both tasks in every corpus. Supplementary Table~S1 reports the point estimate and cluster interval for each encoder and task.
\paperfigure{fig_ranking}{Native logistic-probe CP-AUC and FITB accuracy on the clean test partitions. Each panel orders encoders within its own corpus. Intervals reflect the dependence structure of that corpus.}{fig:ranking}

Validation selects SigLIP2-So400m as the generic encoder on all corpora. It selects Marqo-FashionCLIP on Polyvore and Marqo-FashionSigLIP on IQON and POG in the fashion group. On POG the test leader is PE-Core, even though the validation rule selects SigLIP2-So400m. We keep the validation choice for the contextual and external comparisons. Across corpora the primary pairwise tests distinguish \resHolmSigPrimary{} of \resHolmComparisonsPrimary{} CP-AUC contrasts and \resHolmSigPrimaryFitb{} of \resHolmComparisonsPrimaryFitb{} FITB contrasts after adjustment. The rank intervals in Fig.~\ref{fig:mainrank} show which adjacent positions remain plausible under resampling.

\subsection{Polyvore: gains in whole-outfit discrimination and completion}
\showPolyvore
On Polyvore, CLIP-L reaches 0.907 CP-AUC and 66.7\% FITB accuracy. Both DINOv3 resolutions score lower on both tasks, and both exceed ResNet-50 (0.744 AUC, 44.8\% FITB). All four candidates for a blank share its category (Table~\ref{tab:Polyvore}).

SigLIP2-So400m and PE-Core are close on both tasks: 0.933 and 0.931 AUC, and 72.3\% and 72.1\% FITB accuracy. Qwen-Embedding-8B narrowly exceeds SigLIP2-B on AUC (0.921 versus 0.919), but recovers fewer held-out items (68.2\% versus 69.4\%).

Polyvore validation selects Marqo-FashionCLIP. On test its AUC is 0.917, just below Marqo-FashionSigLIP at 0.918. We retain the validation choice for the transformer and A100 comparisons.

\subsection{IQON: a multimodal embedding leads whole-outfit prediction}
\showIQON
On IQON, Qwen-Embedding-8B has the highest observed AUC (0.788), while SigLIP2-So400m recovers the most completions (49.6\%). Both Qwen embeddings exceed the two Marqo checkpoints on AUC (Table~\ref{tab:IQON}). Each Qwen embedding supplies one vector per item to the shared classifier.

IQON has 22,831 test outfits, of which 54.2\% lie in a single item component. For Qwen-Embedding-8B the crossed-bootstrap interval extends 0.030 below the observed AUC and 0.013 above it. Resampling users alone gives narrower intervals on the same predictions.

\subsection{POG: compatibility in a commercial catalog}
\showPOG
PE-Core leads on both POG tasks (0.763 AUC and 46.2\% FITB), followed by SigLIP2-So400m (0.759 and 45.2\%). The paired intervals include zero for both differences. ResNet-50 has the lowest scores (0.597 and 31.4\%) on this previously unevaluated corpus (Table~\ref{tab:POG}).

Increasing the Qwen embedding model from 2B to 8B parameters raises POG AUC from 0.744 to 0.756. Completion accuracy changes little (42.8\% to 42.6\%). Both trail SigLIP2-B at 43.5\% on choices within the catalog category of the missing item.

\paperfigure{fig_rank_intervals}{Positions of the native probes under cluster resampling. Adjacent positions at the top of a panel can overlap even when the leading encoders are separated from the supervised CNN. Lower numerical rank is better.}{fig:mainrank}

\subsection{What changes with fashion adaptation?}
Marqo-FashionSigLIP improves on its SigLIP parent in native CP-AUC on every corpus, with the clearest point-estimate gain on IQON (Fig.~\ref{fig:pairs}). FashionCLIP~2 trails its OpenCLIP parent on Polyvore and POG, and the IQON estimates round to the same value. GR-Lite trails its resolution-matched DINOv3 parent in native CP-AUC on all three corpora. Projection changes the result. GR-Lite wins all nine Polyvore PCA comparisons after Holm adjustment, and all three IQON PCA-logistic contrasts favor GR-Lite (Table~\ref{tab:signs}). Its Polyvore completion accuracy nevertheless exceeds its parent's. A change in representation can therefore help discriminate candidates for a fixed context even when it lowers the score for separating observed and reconstructed outfits.
Marqo-FashionSigLIP gains in 74 of 78 cells, including ten significant gains in 13 Polyvore CP settings. FashionCLIP~2 consistently loses on Polyvore.
\begin{table}[htbp]\centering\small\setlength{\tabcolsep}{2pt}
\caption{Direction of the fashion checkpoint minus its parent. A star marks Holm-adjusted $p<0.05$. N denotes native features, S adds the coordinate sum, and numbers denote PCA dimensions. L, X and M denote logistic regression, XGBoost and MLP.}\label{tab:signs}
\begin{adjustbox}{max width=\linewidth}\begin{tabular}{llllllllllllllll}\toprule
Checkpoint & Corpus & Task & NL & NM & NX & SL & 128L & 128M & 128X & 256L & 256M & 256X & 512L & 512M & 512X\\\midrule
FashionCLIP 2 & Polyvore & CP & $-^{*}$ & $-$ & $-^{*}$ & $-^{*}$ & $-^{*}$ & $-^{*}$ & $-$ & $-$ & $-^{*}$ & $-$ & $-$ & $-$ & $-^{*}$\\
FashionCLIP 2 & Polyvore & FITB & $-$ & $-^{*}$ & $-$ & $-$ & $-$ & $-$ & $-$ & $-$ & $-$ & $-$ & $-$ & $-^{*}$ & $-$\\
FashionCLIP 2 & IQON & CP & $+$ & $+$ & $+$ & $-$ & $-$ & $-$ & $-$ & $-$ & $-$ & $-$ & $-$ & $-$ & $-$\\
FashionCLIP 2 & IQON & FITB & $-$ & $+$ & $-$ & $+$ & $-$ & $-$ & $-$ & $-$ & $+$ & $-$ & $-$ & $-$ & $-$\\
FashionCLIP 2 & POG & CP & $-$ & $-$ & $+$ & $-$ & $+$ & $-$ & $+$ & $-$ & $+$ & $-$ & $-$ & $-$ & $-$\\
FashionCLIP 2 & POG & FITB & $-$ & $-$ & $-$ & $-$ & $+$ & $+$ & $-$ & $-$ & $-$ & $+$ & $+$ & $-$ & $-$\\
Marqo-FashionSigLIP & Polyvore & CP & $+$ & $+^{*}$ & $+$ & $+$ & $+^{*}$ & $+^{*}$ & $+^{*}$ & $+^{*}$ & $+^{*}$ & $+^{*}$ & $+^{*}$ & $+^{*}$ & $+^{*}$\\
Marqo-FashionSigLIP & Polyvore & FITB & $+$ & $-^{*}$ & $+$ & $+$ & $+^{*}$ & $+$ & $+^{*}$ & $+^{*}$ & $+^{*}$ & $+^{*}$ & $+$ & $+^{*}$ & $+^{*}$\\
Marqo-FashionSigLIP & IQON & CP & $+^{*}$ & $+$ & $+$ & $+^{*}$ & $+^{*}$ & $+$ & $+^{*}$ & $+$ & $+$ & $+^{*}$ & $+$ & $+$ & $+^{*}$\\
Marqo-FashionSigLIP & IQON & FITB & $+$ & $+$ & $+$ & $+$ & $+$ & $+$ & $+$ & $+$ & $+$ & $+$ & $+$ & $+$ & $+$\\
Marqo-FashionSigLIP & POG & CP & $+$ & $+$ & $+$ & $+$ & $+$ & $-$ & $+$ & $+$ & $+$ & $+$ & $+$ & $+$ & $+$\\
Marqo-FashionSigLIP & POG & FITB & $+$ & $+$ & $+$ & $+$ & $+$ & $-$ & $-$ & $+$ & $+$ & $+$ & $+$ & $+$ & $+$\\
GR-Lite & Polyvore & CP & $-$ & $-^{*}$ & $-^{*}$ & $-$ & $+^{*}$ & $+^{*}$ & $+^{*}$ & $+^{*}$ & $+^{*}$ & $+^{*}$ & $+^{*}$ & $+^{*}$ & $+^{*}$\\
GR-Lite & Polyvore & FITB & $+^{*}$ & $+$ & $+$ & $+^{*}$ & $+^{*}$ & $+^{*}$ & $+^{*}$ & $+^{*}$ & $+^{*}$ & $+^{*}$ & $+^{*}$ & $+^{*}$ & $+^{*}$\\
GR-Lite & IQON & CP & $-^{*}$ & $-^{*}$ & $-^{*}$ & $-^{*}$ & $+$ & $-$ & $-$ & $+$ & $-$ & $-$ & $+$ & $-$ & $-$\\
GR-Lite & IQON & FITB & $-$ & $-^{*}$ & $-$ & $-$ & $+$ & $+$ & $+$ & $+$ & $+$ & $+$ & $+$ & $+$ & $+$\\
GR-Lite & POG & CP & $-$ & $-^{*}$ & $-^{*}$ & $-^{*}$ & $-$ & $-$ & $-$ & $-$ & $-$ & $-$ & $-$ & $-$ & $-$\\
GR-Lite & POG & FITB & $-$ & $-$ & $-$ & $-$ & $-$ & $+$ & $-$ & $-$ & $+$ & $-$ & $-$ & $-$ & $-$\\
\bottomrule\end{tabular}\end{adjustbox}\end{table}

\paperfigure{fig_matched_pairs}{Fashion-adapted checkpoints compared with the corresponding parent. The sign of a within-pair contrast identifies the direction of transfer for that adaptation.}{fig:pairs}

\showContrasts
In the native setting, FashionCLIP~2 loses 0.0033 AUC on Polyvore and GR-Lite gains 1.51 percentage points on FITB (Table~\ref{tab:Contrasts}). On IQON, Marqo-FashionSigLIP gains 0.0089 AUC and GR-Lite loses 0.0422. On POG the GR-Lite estimate is similar ($-0.0419$), but the adjusted probability is 0.084.

On IQON, Marqo-FashionSigLIP improves on the parent from which it was adapted. Its native AUC of 0.753 still falls below the validation-selected SigLIP2-So400m at 0.784.

\subsection{Which conclusions survive changes in measurement?}
Across the 39 corpus--readout blocks, compatibility rankings agree closely (Friedman $\chi^2_{14}=458.5$, $p<10^{-88}$, Kendall $W=0.84$). Mean ranks range from PE-Core (1.28), SigLIP2-So400m (2.31) and Marqo-FashionSigLIP (4.28) to ResNet-50 (14.95). The Nemenyi critical difference is 3.43. These statistics describe repeated measurements on shared questions.

\paperfigure{fig_critical_difference}{Mean CP-AUC ranks across 39 blocks. The Nemenyi critical difference is 3.43 rank positions.}{fig:cd}

Mean CP-AUC across encoders increases with outfit length in all three corpora. On IQON it rises from 0.664 for two-item outfits to 0.845 for outfits with at least six items. Reconstructing a longer outfit disrupts more pairs. The length groups also contain different outfits.

After the five uncertain-provenance checkpoints are removed, SigLIP2-So400m leads native CP-AUC on Polyvore (0.933) and IQON (0.784). PE-Core is second on both and first on POG (0.763), followed there by SigLIP2-So400m (0.759). Supplementary Table~S4 gives the restricted rankings.

Supplementary Section~2.1 compares the rebuilt pilot partitions with the frozen splits (Table~S5).

\paragraph{Changing the projection and the learner}
Supplementary Fig.~S1 compares the 12 alternative readouts with the native logistic probe. On IQON, native nonlinear readouts improve CP-AUC for CLIP-L and both DINOv3 resolutions. The gains do not consistently survive PCA.

\paragraph{Labels and alternative question sets}
With 10\% of Polyvore training outfits, SigLIP2-So400m reaches 0.922 AUC, above the full-data scores of both Marqo encoders. With all training outfits its AUC reaches 0.933. The selected generic encoder improves monotonically with the training fraction on all three corpora (Fig.~\ref{fig:mainlabels}).
\paperfigure{fig_label_fractions}{CP-AUC as supervision for the probe varies. The x-axis measures the fraction of the source training outfits. Validation and test partitions retain their roles at every fraction.}{fig:mainlabels}

Regenerating the questions gives Polyvore CP-AUC between 0.9322 and 0.9335 for SigLIP2-So400m. Its construction-seed SD is 0.0005, compared with 0.0016 when sampling variation is included (Supplementary Fig.~S2). The seed SD is likewise smaller than total SD on IQON and POG (Table~\ref{tab:Robust}).
\showRobust

\paragraph{Whole-outfit exclusion}
The looser audit removes a whole outfit if any item matches a training or validation image by pHash or exceeds 0.95 embedding cosine similarity. It excludes 8,365 of 15,142 Polyvore outfits (55\%), 17,407 of 22,831 IQON outfits (76\%) and 10,518 of 14,692 POG outfits (72\%). The pHash criterion flags 2,471, 9,400 and 3,096 outfits, respectively. On Polyvore the embedding criterion flags 7,605. Supplementary Table~S3 gives counts for each criterion and their union. Similar products photographed against white backgrounds can trigger the audit. Each filtered estimate therefore describes a minority of the original outfits with a different product mix. We use it to assess sensitivity to that change in composition.

On the retained outfits, SigLIP2-So400m leads all three corpora on CP-AUC. Its IQON AUC falls from 0.784 to 0.760, while FITB rises from 49.6\% to 51.1\%. These changes reflect which test outfits remain. The classifiers keep their original fits.

\subsection{Do the findings transfer to other outfit decisions?}
\paragraph{A learned outfit scorer}
When OutfitTransformer receives the selected generic features, mean test CP-AUC reaches \resOtPolyvoreVTwoCleanGenericTestCpAucMean{} on Polyvore, \resOtIqonThreeZeroZeroZeroVTwoCleanGenericTestCpAucMean{} on IQON and \resOtPogVTwoCleanGenericTestCpAucMean{} on POG across training seeds. With the selected fashion features the corresponding means are \resOtPolyvoreVTwoCleanFashionTestCpAucMean{}, \resOtIqonThreeZeroZeroZeroVTwoCleanFashionTestCpAucMean{} and \resOtPogVTwoCleanFashionTestCpAucMean{}. ResNet-50 has the lowest transformer AUC in each corpus. Contextual learning raises CP-AUC over the corresponding linear probes across all these inputs. On Polyvore, however, the selected generic transformer's mean FITB accuracy of \resOtPolyvoreVTwoCleanGenericTestFitbAccMean{}\% falls below the \resPolyvoreVTwoCleanSiglipTwoSoFourZeroZeromThreeEightFourNativeLogisticTestFitbAcc{}\% of its linear probe. 

After averaging the five seed scores for each outfit, the generic-minus-fashion AUC differences were 0.0088 [0.0071, 0.0105] on Polyvore, 0.0116 [0.0084, 0.0165] on IQON and 0.0191 [0.0141, 0.0325] on POG (2,000 paired cluster draws, percentile 95\% intervals). All three intervals exclude zero.
\showContext

For completion, the generic transformer improves on its linear probe on IQON (57.3\% versus 49.6\%) and POG (50.1\% versus 45.2\%). Both selected inputs lose completion accuracy on Polyvore. Early stopping used whole-outfit AUC.

\paragraph{External judgments}
Among 15 encoders, the Spearman correlations between native source CP-AUC and expert accuracy were 0.86 [0.61, 0.92] for Polyvore, 0.86 [0.69, 0.91] for IQON and 0.76 [0.33, 0.83] for POG. For crowd-majority agreement the corresponding correlations were 0.29 [0.03, 0.46], 0.53 [0.16, 0.67] and 0.38 [$-0.01$, 0.61]. The 95\% intervals resample the 100 questions 2,000 times. The two marginal intervals are disjoint for Polyvore and IQON and overlap for POG. SigLIP2-So400m recovers 43--51\% of majority choices across source fits, compared with 58--73\% for Qwen-Embedding-8B. The selected fashion probe exceeds SigLIP2-So400m on each source, but only the Polyvore comparison survives adjustment (70 versus 50 correct, $p_{\rm Holm}=0.00525$). The broader roster gives no basis for attributing this difference to fashion training alone.

\showExternal

On the 100 expert questions, both IQON-selected probes score 64\%. The Polyvore generic probe leads its fashion counterpart by one correct answer (73 versus 72). With POG training, the fashion probe leads by eight (72 versus 64). Supplementary Table~S2 breaks these scores down by aesthetic facet. All source fits face the same questions.

On LAT, Qwen-Instruct matches 75 archive answers and 70 crowd majorities. Its choices receive 56.4\% of votes on average, including votes for alternatives that lose the majority.

\paragraph{Zero-shot choice}
On four-option questions the instruction model picked A in 14--17\% of cases, below the 25\% expected under uniform choice. Permuting options across questions distributes the position effect across candidate items. Table~\ref{tab:Paired} compares source-trained probes on the same questions. Table~\ref{tab:External} gives A100 zero-shot performance with question-bootstrap intervals.
\paperfigure{fig_qualitative_fitb}{Each corpus contributes examples recovered by its validation-selected generic probe and an example missed by the entire probe roster. The marked choices refer to the generic probe. The target is the item present in the recorded outfit.}{fig:cases}

All 15 probes miss the illustrated Polyvore target: the selected wallet and recorded bag both match the warm orange tones of the outfit. In the IQON example the selected earrings repeat the citrus-slice motif of the bag (Fig.~\ref{fig:cases}).

\showPaired
On common Polyvore FITB questions the instruction model scores 57.0\%, between ResNet-50 (44.4\%) and the selected fashion probe (65.4\%). On IQON the generic probe leads zero-shot choice by 10.1 percentage points ($p_{\rm Holm}=0.018$). The POG difference is 7.3 points ($p_{\rm Holm}=0.424$, Table~\ref{tab:Paired}).

\paragraph{Retrieval ranks}
Only six checkpoints overlap with LookBench. Their rank correlations range from $-0.14$ to $0.14$, with 95\% intervals extending roughly from $-0.86$ to $0.84$. This panel leaves both the sign and magnitude of the association unresolved.

\showRetrieval
GR-Lite has the best retrieval score in the shared panel (65.71) and the lowest native CP-AUC point estimate among the vision transformers on all three corpora. On Polyvore, its AUC of 0.847 is close to the 0.849 of resolution-matched DINOv3 ($p_{\rm Holm}=1.0$). CLIP-L, last among the six on retrieval (39.79), exceeds GR-Lite on each compatibility test (Table~\ref{tab:Retrieval}).

\subsection{Calibration and extraction cost}

\showCalibrated
On Polyvore, Platt scaling reduces the Brier score of SigLIP2-So400m from 0.1036 to 0.0991, while log loss increases from 0.3501 to 0.3516. On IQON both mappings reduce log loss but increase 15-bin ECE. On POG, Platt scaling reduces ECE from 0.0139 to 0.0069 (Table~\ref{tab:Calibrated}). Each test contains equal numbers of observed and generated outfits.

\paperfigure{fig_calibration}{Held-out calibration error before and after fitting a score map on source validation examples. The same trained classifier supplies each set of scores.}{fig:maincal}

\paragraph{An encoder for a growing catalog}
\showCost
ResNet-50 has the lowest single-image latency (2.37 ms). With batches of 64, OpenCLIP-B/32 has the highest throughput (4,493 images/s). It also exceeds ResNet-50 on Polyvore FITB, 66.7\% to 44.8\% (Table~\ref{tab:Cost}).

SigLIP-B and Marqo-FashionSigLIP each take 4.73 ms per image at batch size one. The adaptation gains 0.0089 native AUC on IQON at the same measured latency. SigLIP2-So400m reaches 0.784 AUC there and takes 13.73 ms.

Qwen-Embedding-8B peaks at 17.47 GB, compared with 3.37 GB for SigLIP2-So400m. The Qwen parameter count covers the vision tower, while latency, memory and energy include the language stack that produces the final vector. Supplementary Table~S6 adds operation counts, batch-64 latency, CPU preprocessing and CO$_2$e.
